\documentclass[a4paper, 12pt, twoside]{article}
\usepackage[utf8]{inputenc}
\usepackage[T1]{fontenc}		
\usepackage[francais]{babel}
\usepackage{lmodern}
\usepackage{ae,aecompl}
\usepackage[top=2.5cm, bottom=2cm, 
			left=3cm, right=2.5cm,
			headheight=15pt]{geometry}
\usepackage{graphicx}
\usepackage{eso-pic}
\usepackage{array} 
\usepackage{hyperref}
\usepackage{xcolor}
\usepackage{listings}
\usepackage{enumitem}

% ── Style du code ──────────────────────────────────────────
\definecolor{codegris}{RGB}{245,245,245}
\definecolor{codevert}{RGB}{0,128,0}
\definecolor{codebleu}{RGB}{0,0,200}
\lstset{
  backgroundcolor=\color{codegris},
  basicstyle=\ttfamily\small,
  keywordstyle=\color{codebleu}\bfseries,
  commentstyle=\color{codevert},
  frame=single,
  breaklines=true,
  showstringspaces=false,
  numbers=left,
  numberstyle=\tiny\color{gray},
  language=Python,
}

\input{pagedegarde}

\title{Détection automatique des Fake News}
\entreprise{Université Paris Nanterre}
\datedebut{1 Avril 2025}
\datefin{23 Mai 2025}

\membrea{\textbf{OUATIK} \textbf{Marwane} \textcolor{red}{\textbf{45018118}}}
\membreb{\textbf{GHANEM} \textbf{Rayan} \textcolor{red}{\textbf{45017436}}}
\membrec{\textbf{BENDJELLOUL} \textbf{Rayane} \textcolor{red}{\textbf{4501665}}}
\membred{\textbf{BENELKHAZNADJI} \textbf{Nadir} \textcolor{red}{\textbf{4501665}}}

\begin{document}
\pagedegarde

\tableofcontents
\newpage

% ============================================================
\section{Introduction}
% ============================================================

Dans un monde où l'information circule à une vitesse sans précédent, la
désinformation constitue un danger croissant pour les sociétés démocratiques.
Les \textit{fake news} --- ou fausses informations --- se propagent via les
réseaux sociaux, les médias en ligne et les messageries instantanées, rendant
difficile pour le grand public de distinguer le vrai du faux.

Dans ce projet, nous avons conçu un système automatique de détection des
fake news s'appuyant sur des techniques d'apprentissage automatique
(\textit{Machine Learning}). Pour ce faire, nous utilisons deux jeux de
données : \texttt{Fake.csv} et \texttt{True.csv}, contenant respectivement
des articles faux et des articles authentiques.

Ce rapport décrit les choix techniques effectués, les outils utilisés, les
difficultés rencontrées et les résultats obtenus.

% ============================================================
\section{Environnement de travail}
% ============================================================

Nous avons travaillé en \textbf{Python}, langage incontournable en Data
Science grâce à sa syntaxe claire et sa richesse en bibliothèques spécialisées.

Concernant les environnements de développement, nous avons utilisé deux outils
complémentaires :

\begin{itemize}
  \item \textbf{JupyterLite (WASM)} : environnement principal d'exécution,
        accessible directement via un navigateur web sans installation locale.
        Les bibliothèques s'installent via la commande \texttt{!mamba install}.
  \item \textbf{Visual Studio Code} : utilisé pour l'édition locale du code
        source, grâce à sa légèreté et ses nombreuses extensions Python.
\end{itemize}

Pour la rédaction du rapport, nous avons utilisé \textbf{Overleaf}, une
plateforme en ligne de rédaction \LaTeX{} permettant la collaboration en
temps réel.

% ============================================================
\section{Description du projet et objectifs}
% ============================================================

\subsection{Mise en contexte}

Le projet s'appuie sur deux fichiers CSV issus d'un dataset public largement
utilisé dans la recherche sur la désinformation. \texttt{Fake.csv} contient
des articles de presse identifiés comme faux, tandis que \texttt{True.csv}
contient des articles authentiques. Après fusion et mélange aléatoire, le
dataset comprend plusieurs dizaines de milliers d'articles étiquetés.

L'idée centrale est d'entraîner un modèle de classification supervisée
capable, à partir du texte brut d'un article, de prédire s'il s'agit d'une
fake news ou d'une information authentique.

\subsection{Objectifs}

Les objectifs de notre projet sont :

\begin{itemize}
  \item \textbf{Traitement du langage naturel (NLP)} : nettoyer et préparer
        les textes (suppression des stopwords, URLs, caractères spéciaux).
  \item \textbf{Classification binaire} : distinguer les fausses informations
        des vraies avec une précision maximale.
  \item \textbf{Détection mathématique} : identifier et vérifier
        automatiquement les expressions mathématiques saisies par l'utilisateur.
  \item \textbf{Interface utilisateur} : permettre à un utilisateur de
        soumettre un texte et d'obtenir un verdict
        \texttt{FAKE NEWS} ou \texttt{AUTHENTIC NEWS}.
  \item \textbf{Sauvegarde du modèle} : exporter le modèle entraîné pour
        une réutilisation sans ré-entraînement.
\end{itemize}

% ============================================================
\section{Bibliothèques, Outils et technologies}
% ============================================================

\subsection*{Bibliothèques Python utilisées}

\begin{itemize}
  \item \texttt{\textbf{pandas}} : chargement et manipulation des datasets CSV.
  \item \texttt{\textbf{numpy}} : opérations numériques de base.
  \item \texttt{\textbf{scikit-learn}} : vectorisation TF-IDF
        (\texttt{TfidfVectorizer}), modèle \texttt{LogisticRegression},
        métriques d'évaluation (\texttt{accuracy\_score},
        \texttt{classification\_report}, \texttt{confusion\_matrix}).
  \item \texttt{\textbf{nltk}} : traitement du langage naturel, suppression
        des stopwords anglais.
  \item \texttt{\textbf{joblib}} : sérialisation et sauvegarde du modèle
        et du vectoriseur entraînés.
  \item \texttt{\textbf{ast / re / operator}} : analyse syntaxique et
        évaluation sécurisée des expressions mathématiques (sans
        \texttt{eval()}).
\end{itemize}

\subsection*{Technologies clés}

\begin{itemize}
  \item \textbf{TF-IDF} (\textit{Term Frequency--Inverse Document Frequency}) :
        méthode de vectorisation qui transforme les textes en vecteurs
        numériques en pondérant l'importance relative de chaque mot.
  \item \textbf{Régression Logistique} : algorithme de classification binaire
        supervisée, simple, interprétable et très performant sur des
        représentations TF-IDF.
\end{itemize}

% ============================================================
\section{Travail réalisé}
% ============================================================

\subsection*{Répartition des tâches}

Nous avons divisé le travail en trois parties afin que chacun puisse
contribuer pleinement au projet :

\begin{itemize}
  \item \textbf{Partie 1 --- OUATIK Marwane} : Préparation des données,
        chargement des CSV, nettoyage des textes (fonction
        \texttt{clean\_text}), étiquetage et fusion des datasets.
  \item \textbf{Partie 2 --- GHANEM Rayan} : Vectorisation TF-IDF,
        entraînement du modèle de régression logistique, évaluation
        (accuracy, matrice de confusion, rapport de classification),
        sauvegarde du modèle avec \texttt{joblib}.
  \item \textbf{Partie 3 --- BENDJELLOUL Rayane et BENELKHAZNADJI Nadir} :
        Détection et évaluation sécurisée des expressions mathématiques
        via le module \texttt{ast}, boucle interactive utilisateur.
\end{itemize}

\subsection*{Fonctionnalités prévues et réalisées}

\begin{table}[h]
\begin{center}
\begin{tabular}{|p{7cm}|c|p{4cm}|}
\hline
\textbf{Fonctionnalité} & \textbf{Réalisée ?} & \textbf{Remarques} \\
\hline
Chargement et fusion des datasets & Oui & --- \\
\hline
Nettoyage NLP des textes & Oui & --- \\
\hline
Vectorisation TF-IDF & Oui & --- \\
\hline
Entraînement régression logistique & Oui & Précision $>$ 98\,\% \\
\hline
Sauvegarde du modèle (.pkl) & Oui & --- \\
\hline
Détection mathématique sécurisée & Oui & Sans \texttt{eval()} \\
\hline
Interface utilisateur interactive & Oui & Boucle \texttt{while} \\
\hline
Interface web (Flask/Streamlit) & Non & Manque de temps \\
\hline
\end{tabular}
\end{center}
\caption{Liste des fonctionnalités prévues et leur état de réalisation}
\label{tab:fonctionnalites}
\end{table}

Le tableau \ref{tab:fonctionnalites} résume l'ensemble des fonctionnalités.
L'interface web n'a pas pu être réalisée par manque de temps lors de la
phase finale du projet.

% ============================================================
\section{Difficultés rencontrées}
% ============================================================

\begin{itemize}
  \item \textbf{Compréhension des données} : Il a fallu analyser le contenu
        des fichiers CSV et s'assurer que les labels étaient correctement
        attribués avant tout entraînement.

  \item \textbf{Environnement JupyterLite (WASM)} : L'installation des
        bibliothèques nécessite \texttt{mamba} et non \texttt{pip}, ce qui
        est inhabituel. La gestion des chemins de fichiers dans un
        environnement navigateur a également posé des problèmes.

  \item \textbf{Sécurité de l'évaluateur mathématique} : L'utilisation de
        \texttt{eval()} étant dangereuse, nous avons dû implémenter un
        évaluateur basé sur le module \texttt{ast}, qui analyse l'arbre
        syntaxique abstrait de l'expression.

  \item \textbf{Limites des LLM} : Nous avons utilisé ChatGPT et Mistral
        comme outils d'assistance. Si utiles pour déboguer, ils ont parfois
        produit du code qui semblait correct mais ne fonctionnait pas dans
        notre environnement spécifique.
\end{itemize}

\newpage
% ============================================================
\section{Bilan}
% ============================================================

\subsection{Conclusion}

Ce projet nous a permis de mettre en pratique les grandes étapes d'un
pipeline de Machine Learning appliqué au traitement du langage naturel :
préparation des données, extraction de features, entraînement, évaluation
et déploiement sous forme d'interface interactive. Le modèle atteint une
précision supérieure à 98\,\% sur les données de test, ce qui confirme
l'efficacité de l'approche TF-IDF + régression logistique pour ce type
de tâche.

\subsection{Perspectives}

Plusieurs améliorations pourraient être apportées à notre système :

\begin{itemize}
  \item Tester des modèles plus sophistiqués : transformers (BERT, RoBERTa)
        pour une meilleure compréhension sémantique.
  \item Analyser également le titre et la source de l'article.
  \item Créer une interface web avec Flask ou Streamlit.
  \item Étendre la détection à d'autres langues, notamment le français.
\end{itemize}

% ============================================================
\section{Webographie}
% ============================================================
\begin{thebibliography}{4}
  \bibitem[SKL-DOC]{skldoc} \url{https://scikit-learn.org/stable/}
  \bibitem[NLTK]{nltk} \url{https://www.nltk.org/}
  \bibitem[DATASET]{dataset} \url{https://www.kaggle.com/clmentbisaillon/fake-and-real-news-dataset}
  \bibitem[OVERLEAF]{overleaf} \url{https://www.overleaf.com/}
  \bibitem[GPT]{gpt} \url{https://chat.openai.com/}
  \bibitem[MISTRAL]{mistral} \url{https://chat.mistral.ai/}
\end{thebibliography}

\newpage

% ============================================================
\section{Annexes}
% ============================================================
\appendix
\makeatletter
\def\@seccntformat#1{Annexe~\csname the#1\endcsname:\quad}
\makeatother

% ── ANNEXE A ────────────────────────────────────────────────
\section{Cahier des charges}

\textbf{1. Introduction}\\
Le projet vise à développer un programme de détection automatique des fake
news à partir de textes en anglais, en s'appuyant sur le NLP et le Machine
Learning.

\medskip
\textbf{2. Objectifs}
\begin{itemize}
  \item Développer un pipeline complet de traitement et classification de textes.
  \item Offrir une interface interactive pour soumettre un article.
  \item Gérer la vérification des expressions mathématiques.
  \item Sauvegarder le modèle entraîné pour réutilisation.
\end{itemize}

\medskip
\textbf{3. Fonctionnalités}
\begin{itemize}
  \item Nettoyage NLP : stopwords, URLs, caractères spéciaux.
  \item Vectorisation TF-IDF.
  \item Classification binaire : FAKE / AUTHENTIC.
  \item Évaluation mathématique sécurisée (sans \texttt{eval()}).
\end{itemize}

\medskip
\textbf{4. Planning : du 1 Avril au 23 Mai 2025}
\begin{itemize}
  \item Phase 1 : Analyse des besoins (1 semaine).
  \item Phase 2 : Préparation des données et pipeline NLP (2 semaines).
  \item Phase 3 : Entraînement et évaluation du modèle (1 semaine).
  \item Phase 4 : Interface utilisateur et tests finaux (1 semaine).
\end{itemize}

% ── ANNEXE B ────────────────────────────────────────────────
\section{Exemple d'exécution du projet}

Pour cet exemple, nous avons testé les entrées suivantes dans la boucle
interactive :

\begin{lstlisting}[language={}, title={Sortie console}]
================================
DETECTION AUTOMATIQUE DES FAKE NEWS GRACE A L'IA
================================

Entree : 2 + 2 = 4
Resultat : AUTHENTIC NEWS

Entree : 10 > 5
Resultat : AUTHENTIC NEWS

Entree : 5 * 7 = 20
Resultat : FAKE NEWS

Entree : NASA discovered water on Mars
Resultat : AUTHENTIC NEWS

Entrez une information (ou tapez exit) : exit
Programme arrete
\end{lstlisting}

Les trois premiers exemples illustrent la détection mathématique :
$2+2=4$ est vrai (AUTHENTIC), $10>5$ est vrai (AUTHENTIC), mais
$5 \times 7 = 20$ est faux (FAKE NEWS). Le dernier exemple illustre
la classification NLP sur un texte journalistique.

% ── ANNEXE C ────────────────────────────────────────────────
\section{Manuel utilisateur}

\textbf{Installation des dépendances}

Dans un environnement JupyterLite :
\begin{lstlisting}[language=bash, numbers=none]
!mamba install pandas numpy scikit-learn nltk joblib
\end{lstlisting}

Dans un environnement Python standard :
\begin{lstlisting}[language=bash, numbers=none]
pip install pandas numpy scikit-learn nltk joblib
\end{lstlisting}

\medskip
\textbf{Fichiers nécessaires}
\begin{itemize}
  \item \texttt{Fake.csv} --- dataset des articles faux
  \item \texttt{True.csv} --- dataset des articles authentiques
  \item Le notebook ou script Python principal
\end{itemize}

\medskip
\textbf{Structure du dossier recommandée}
\begin{lstlisting}[language=bash, numbers=none]
projet_fakenews/
  |-- fake_news_detection.ipynb
  |-- Fake.csv
  |-- True.csv
  |-- fake_news_model.pkl      (genere apres entrainement)
  |-- tfidf_vectorizer.pkl     (genere apres entrainement)
\end{lstlisting}

\medskip
\textbf{Exécution}
\begin{enumerate}
  \item Placer \texttt{Fake.csv} et \texttt{True.csv} dans le même dossier
        que le script.
  \item Exécuter toutes les cellules dans l'ordre.
  \item Saisir un texte ou une expression mathématique à l'invite.
  \item Taper \texttt{exit} pour quitter le programme.
\end{enumerate}

\end{document}
